\chapter{Integrals that collapse}

\begin{orient}
\textbf{What this chapter is for.} This is the mirror image of the chapter on special functions. There, an integral that looked routine turned out to require the dilogarithm or the Gamma function. Here, an integral that looks impossible turns out to require nothing at all: the integrand is exactly the derivative of something you can write down, or it telescopes, or it is a constant wearing a costume, or an odd factor on a symmetric interval annihilates it, or the only work is a case split on an absolute value. In every case the entire difficulty is noticing, and once you have noticed there is no integration to do. The species are few enough to enumerate, and the enumeration is the point of the chapter: perfect derivatives (of a product, of a quotient, of a power tower), telescoping sums and shells, disguised constants and disguised zeros, step functions that only need their level sets, and the one genuinely dangerous species, an integrand whose naive antiderivative is discontinuous inside the interval and therefore wrong. The chapter ends with a diagnostic sweep that catches all of them in under a minute, and the sweep is worth more than any individual technique in it, because these problems are designed so that the only alternative to noticing is an hour of futile work.
\end{orient}

\section{The sweep}

The problems in this chapter are adversarial. They are written so that a competent integrator, reading left to right, will start integrating; and the integration will not terminate, because there is nothing there to integrate. The defence is a fixed sequence of cheap tests applied \emph{before} any method is chosen. Each test costs a few seconds, and each one is a question about structure rather than about technique.

\begin{keynote}[The one-minute sweep]
\textbf{1. Evaluate the integrand at two points.} Equal values on a non-constant-looking expression means a disguised constant. Cost: two arithmetic evaluations.\\
\textbf{2. Test parity about the midpoint.} On $[-a,a]$, or after $t=x-\tfrac{a+b}{2}$: is any factor odd against an even remainder?\\
\textbf{3. Look for a quotient with a squared denominator} and a numerator of unexpectedly low degree. That is $\bigl(u/v\bigr)'$.\\
\textbf{4. Look for a sum of unrelated pieces} where one factor is a power tower or a stray logarithm. That is $F'$ found by logarithmic differentiation.\\
\textbf{5. Check whether the interval decomposes.} Floors, fractional parts, absolute values, $\max$, $\min$, $\sgn$: the integrand is smooth on each level set.\\
\textbf{6. Check the domain and the branch.} Even roots, odd roots, $\arcsec$, $\arctan$ of something that crosses a pole, and every denominator zero inside the interval.\\
If all six fail, the integral is genuine and belongs to another chapter.
\end{keynote}

The sweep is ordered by cost, not by frequency. Step 1 is first because it is nearly free and because its payoff is total: a disguised constant reduces the problem to multiplication. Step 6 is last because it requires the most thought, and because it is the only step that catches an error rather than an opportunity.

It helps to understand how these problems are manufactured, because the construction is always the same and always leaves fingerprints. The author starts with the answer, not with the question. They pick an $F$, differentiate it, and then obscure the result: expand a product so that the two halves of the product rule land in separate terms; write $x^{1/x-2}$ instead of $\tfrac{x^{1/x}}{x^{2}}$ so the power tower is invisible; insert a decorative factor that is odd on a symmetric interval; wrap a constant in an identity that hides it. Each of those moves leaves a residue. A sum of terms with unrelated shapes is the fingerprint of a product rule expanded. An exponent written as a difference is the fingerprint of a quotient. A high odd power multiplied by an even weight is the fingerprint of a parity kill. A parameter that appears in the problem and could not possibly appear in the answer is the fingerprint of a cancellation. Reading fingerprints is faster than reading integrands, and it is what the sweep operationalises.

One warning about the sweep before the techniques begin. Its purpose is to \emph{detect}, not to prove. A sample point that matches is evidence for a disguised constant, and a squared denominator is evidence for a quotient rule, but every one of the six tests must be followed by a symbolic confirmation: differentiate the guessed $F$ and compare, or establish the identity that makes the integrand constant. Detection is cheap and unreliable; confirmation is cheap and reliable; and the whole value of the sweep lies in doing the first so that you know which of the second to attempt.

\section{Derivatives wearing a costume}

The largest species by far is the integrand that is literally $F'$ for a writable $F$. What makes these hard is that the standard differentiation rules are being read backwards, and the product rule read backwards does not look like a product: it looks like a sum of two unrelated terms. The two techniques below are the two rules most often disguised, and they are worth separating because their cues are different.

\tech{Perfect-derivative disguise}{core}
\cue An integrand that is a \emph{sum} of pieces which do not obviously belong together, especially when one of them is a power tower $f^{g}$ such as $x^{1/x}$, $x^{x}$, $(\sec x)^{\sin x}$, or when a stray logarithm appears in a product where nothing else suggests one.
\method Guess $F$ from the most complicated factor. Compute $F'$ by logarithmic differentiation,
\[F'=F\cdot(\ln F)',\qquad\text{so for }F=u^{v}:\quad F'=u^{v}\Bigl(v'\ln u+\frac{vu'}{u}\Bigr),\]
and compare with the whole integrand. If it matches, the answer is $F(b)-F(a)$ and no integration happens at all. The tell for a power tower is that its derivative always carries the factor $u^{v}$ itself, so the integrand contains a copy of the very expression you are trying to integrate.

\begin{worked}
\[\int_0^{2}x^{\frac1x-2}\bigl(1-\ln x\bigr)\dd x.\]
The exponent $\tfrac1x-2$ is $\tfrac1x$ minus $2$, which suggests $F=x^{1/x}$ divided by $x^{2}$. Confirm: $\ln F=\dfrac{\ln x}{x}$, so $(\ln F)'=\dfrac{1-\ln x}{x^{2}}$ and
\[F'=x^{1/x}\cdot\frac{1-\ln x}{x^{2}}=x^{\frac1x-2}(1-\ln x),\]
which is the integrand exactly. Hence $I=\bigl[x^{1/x}\bigr]_0^{2}=2^{1/2}-0=\sqrt2\approx1.414214$. Quadrature returns $1.4142135624$.
\end{worked}

\begin{worked}[The plainer sibling]
\[\int_1^{2}x^{x}(1+\ln x)\dd x.\]
Here $F=x^{x}$ has $\ln F=x\ln x$ and $(\ln F)'=1+\ln x$, so $F'=x^{x}(1+\ln x)$ is the integrand. The value is $\bigl[x^{x}\bigr]_1^{2}=4-1=3$. The pattern is worth memorising in the abstract: whenever you see $u^{v}$ multiplied by something, differentiate $u^{v}$ once and see whether that something is $\bigl(v\ln u\bigr)'$.
\end{worked}

\begin{worked}[A trigonometric power tower]
\[\int_0^{\pi/3}(\sec x)^{\sin x}\bigl(\cos x\ln\sec x+\sin x\tan x\bigr)\dd x.\]
The bracket is unrecognisable until you notice that the leading factor is a power tower. Set $F=(\sec x)^{\sin x}$, so $\ln F=\sin x\ln\sec x$ and
\[(\ln F)'=\cos x\ln\sec x+\sin x\cdot\frac{\sec x\tan x}{\sec x}=\cos x\ln\sec x+\sin x\tan x,\]
which is exactly the bracket. Hence $F'$ is the integrand and
\[I=\bigl[(\sec x)^{\sin x}\bigr]_0^{\pi/3}=2^{\sqrt3/2}-1\approx0.822635,\]
confirmed by quadrature to ten digits. The bracket was constructed by differentiating; the only way to read it is to differentiate the same thing.
\end{worked}

\begin{trap}
The lower endpoint. In the first example $x^{1/x}=\exp\bigl(\tfrac{\ln x}{x}\bigr)\to\exp(-\infty)=0$ as $x\to0^{+}$, because $\tfrac{\ln x}{x}\to-\infty$ faster than any polynomial: at $x=10^{-8}$ the value is already below $10^{-8\times10^{7}}$. That limit must be argued, not assumed, and the integral is genuinely improper at $0$. The second failure mode is matching the guessed $F'$ against only part of the integrand: if $F'$ reproduces one term and leaves a remainder, you have not found a perfect derivative, you have found an integration by parts, and the remainder still has to be integrated.
\end{trap}

\begin{seealsob}Reverse product and quotient rule; logarithmic differentiation; engineering the antiderivative $v$; constant or zero integrand in disguise.\end{seealsob}

Power towers announce themselves. The quotient rule read backwards is quieter, because a quotient with a squared denominator is exactly what partial fractions and trigonometric substitution both produce, so the shape does not by itself suggest that anything special is happening. The diagnostic is arithmetic rather than visual: compare the degree of the numerator against the degree you would expect from an honest partial-fraction problem.

\tech{Reverse product and quotient rule}{core}
\cue An integrand of the form $\dfrac{u'v-uv'}{v^{2}}$, or $f'g+fg'$, or a numerator that looks like a Wronskian $fg'-gf'$ sitting over $f^{2}+g^{2}$.
\method Recognise the three shapes:
\[\dvop{x}\frac{u}{v}=\frac{u'v-uv'}{v^{2}},\qquad \dvop{x}(uv)=u'v+uv',\qquad \dvop{x}\arctan\frac{g}{f}=\frac{fg'-gf'}{f^{2}+g^{2}}.\]
The third is the powerful one, because $f^{2}+g^{2}$ downstairs is a shape that no standard method handles well and that the arctan form handles instantly. Diagnostically: a squared denominator whose numerator has \emph{lower} degree than a partial-fraction problem would produce is the tell for the first shape; a sum of two squares downstairs with a difference of cross terms upstairs is the tell for the third.

\begin{worked}
\[\int_0^{1}\frac{x\eu^{x}}{(1+x)^{2}}\dd x.\]
Partial fractions is useless because of the $\eu^{x}$, and parts produces another integral of the same difficulty. Instead take $u=\eu^{x}$, $v=1+x$:
\[\dvop{x}\frac{\eu^{x}}{1+x}=\frac{\eu^{x}(1+x)-\eu^{x}}{(1+x)^{2}}=\frac{x\eu^{x}}{(1+x)^{2}},\]
the integrand exactly. So $I=\Bigl[\dfrac{\eu^{x}}{1+x}\Bigr]_0^{1}=\dfrac{\eu}{2}-1\approx0.359141$.
\end{worked}

\begin{worked}[The Wronskian shape]
\[\int_0^{\pi/2}\frac{x\cos x-\sin x}{x^{2}+\sin^{2}x}\dd x.\]
The denominator is $f^{2}+g^{2}$ with $f=x$, $g=\sin x$; the numerator is $fg'-gf'=x\cos x-\sin x$. Hence the integrand is $\dvop{x}\arctan\dfrac{\sin x}{x}$ and
\[I=\Bigl[\arctan\frac{\sin x}{x}\Bigr]_0^{\pi/2}=\arctan\frac{2}{\pi}-\arctan1=\arctan\frac{2}{\pi}-\frac{\pi}{4}\approx-0.218487.\]
At the lower limit $\tfrac{\sin x}{x}\to1$, so the endpoint is removable rather than singular. The same pattern with $f=\operatorname{Ai}$, $g=\operatorname{Bi}$ and their constant Wronskian $\tfrac1\pi$ gives $\displaystyle\int\frac{\dd x}{\operatorname{Ai}^{2}+\operatorname{Bi}^{2}}=\pi\arctan\frac{\operatorname{Bi}}{\operatorname{Ai}}$, an integral with no business being elementary.
\end{worked}

\begin{worked}[The product rule, split across two terms]
\[\int_0^{\pi/4}\bigl(x\sec^{2}x+\tan x\bigr)\dd x.\]
The two terms look unrelated: the first invites parts, the second is a standard logarithm. But $\dvop{x}(x\tan x)=\tan x+x\sec^{2}x$ is exactly the sum, so
\[I=\bigl[x\tan x\bigr]_0^{\pi/4}=\frac{\pi}{4}\cdot1-0=\frac{\pi}{4}\approx0.785398.\]
The diagnostic that finds this in five seconds: whenever an integrand is a sum of two terms and one of them is $u$ times the derivative of $v$ while the other is $v$ times the derivative of $u$, the answer is $\bigl[uv\bigr]_a^{b}$. Doing it the long way produces $\bigl[x\tan x\bigr]-\int\tan x$ from parts, and then the second term cancels the leftover, which is the same answer arrived at with three times the writing.
\end{worked}

\begin{trap}
Attacking with partial fractions or parts and never testing the structure. The cost of the test is one differentiation of the obvious candidate $\tfrac{u}{v}$, which takes ten seconds; the cost of missing it is unbounded. The second trap is a sign: $\dvop{x}\dfrac uv$ has $u'v-uv'$ upstairs, and the reversed order $uv'-u'v$ integrates to $-\dfrac uv$. Since both forms look equally natural on the page, check which one you actually have before writing the answer.
\end{trap}

\begin{seealsob}Perfect-derivative disguise; integration by parts, cyclic; engineering the antiderivative $v$; Leibniz rule with variable limits.\end{seealsob}

\section{Nothing there at all}

Step 1 of the sweep catches the most extreme species: an integrand which is not a function of $x$ in any meaningful sense, because every apparent dependence cancels. These problems are built by taking an identity, hiding it, and then asking for an integral, and the identity is usually one of a small stock: $\arctan x+\arctan\tfrac1x=\tfrac{\pi}{2}$ for $x>0$, $\log_ab\cdot\log_bc=\log_ac$, a fixed-point equation for a nested radical, Euler's reflection formula for a product of Gamma values, or a determinant of matched antiderivatives.

\tech{Constant or zero integrand in disguise}{medium}
\cue An intimidating integrand with no free structure: mixed logarithm bases, an infinitely iterated radical, a determinant whose entries are antiderivatives of related functions, a sum of inverse trigonometric functions of reciprocal arguments, a product of shifted Gamma or double-factorial expressions inside a logarithm.
\method Evaluate the integrand at a sample point. Then evaluate it at a second. If the two agree, suspect a constant and prove it: chase the identity rather than the integral. Fixed-point equations collapse iterated radicals ($y=\sqrt{x+y}$ or $y=\sqrt{xy}$); change of base collapses mixed logarithms; the Wronskian-style cancellation collapses determinants; and complementary-angle identities collapse sums of inverse trigonometric functions. If the constant is $0$, look first for an odd factor against an even one on a symmetric interval.

\begin{worked}
\[\int_2^{4}\bigl(\log_x8\bigr)\bigl(\log_2x\bigr)\dd x.\]
Sample at $x=3$: $\log_38\cdot\log_23=\dfrac{\ln8}{\ln3}\cdot\dfrac{\ln3}{\ln2}=\dfrac{\ln8}{\ln2}=3$. Sample at $x=4$: also $3$. The $\ln x$ cancels identically, so the integrand is the constant $3$ and $I=3\cdot2=6$. Nothing was integrated.

The same move on a determinant. With
\[M(x)=\begin{pmatrix}\int_0^{x}\cos t\dd t & \int_0^{x}\sin2t\dd t\\[3pt] \int_0^{x}\dd t & \int_0^{x}(t\cos t+\sin t)\dd t\end{pmatrix}=\begin{pmatrix}\sin x & \sin^{2}x\\ x & x\sin x\end{pmatrix},\]
using $\int_0^{x}\sin2t\dd t=\tfrac{1-\cos2x}{2}=\sin^{2}x$ and $\int_0^{x}(t\cos t+\sin t)\dd t=\bigl[t\sin t\bigr]_0^{x}=x\sin x$, the determinant is $x\sin^{2}x-x\sin^{2}x=0$ for every $x$. So $\int_a^{b}\det M(x)\dd x=0$ on any interval whatever.
\end{worked}

\begin{worked}[A disguised zero by parity, and a disguised $\tfrac{\pi}{2}$]
\[\int_{-2}^{2}\ln\!\Bigl(\frac{3+x}{3-x}\Bigr)\sqrt{4-x^{2}}\dd x=0,\]
because $\ln\tfrac{3+x}{3-x}$ is odd, $\sqrt{4-x^{2}}$ is even, and the product is odd on a symmetric interval; quadrature returns $1.1\times10^{-16}$. The weight was chosen to look like an area computation and the logarithm to look like a partial-fraction problem, and neither has to be touched.

\[\int_0^{1}\Bigl(\arctan x+\arctan\frac1x\Bigr)\dd x=\int_0^{1}\frac{\pi}{2}\dd x=\frac{\pi}{2}\approx1.570796.\]
Both terms separately require integration by parts; their sum is a constant on $(0,\infty)$, and the singularity of $\arctan\tfrac1x$ at $0$ is removable in the sum.
\end{worked}

\begin{trap}
Not checking numerically. One evaluation at $x=1$ costs nothing and settles the question, and without it a reader will spend the whole allotted time integrating a constant. The converse trap: two sample points agreeing is evidence, not proof. $\ln\tfrac{3+x}{3-x}\sqrt{4-x^{2}}$ agrees at $x=1$ and $x=-1$ in absolute value but is not constant, and $\sin(\pi x)$ agrees at every integer. Sample at irrational-looking points, and once you believe the claim, prove the identity.
\end{trap}

\begin{seealsob}Even and odd parity on a symmetric interval; algebraic camouflage collapse; self-similar collapse; numerical cross-check as verification discipline.\end{seealsob}

\begin{keynote}[The stock identities behind a disguised constant]
There are not many, and knowing them turns step 1 of the sweep from a numerical experiment into a recognition. Inverse trigonometric: $\arctan x+\arctan\tfrac1x=\tfrac{\pi}{2}$ for $x>0$ and $-\tfrac{\pi}{2}$ for $x<0$; $\arcsin x+\arccos x=\tfrac{\pi}{2}$; $\arctan a+\arctan b=\arctan\tfrac{a+b}{1-ab}$ modulo $\pi$. Logarithmic: $\log_ab\cdot\log_bc=\log_ac$, so any chain of mixed bases collapses to the endpoints. Gamma: Euler reflection $\Gfun(z)\Gfun(1-z)=\tfrac{\pi}{\sin\pi z}$ turns a product of matched Gamma values into $\pi$ over a sine. Hyperbolic: $\cosh^{2}-\sinh^{2}=1$, and $\operatorname{arsinh}x=\ln(x+\sqrt{x^{2}+1})$, which is how a logarithm of a radical becomes an odd function. Fixed points: $y=\sqrt{xy}$ gives $y=x$, and $y=\sqrt{x+y}$ gives $y=\tfrac{1+\sqrt{1+4x}}{2}$, so an infinite radical is never infinite and is often trivial. Finally, any $2\times2$ determinant whose columns are proportional collapses to $0$, which is what makes matched antiderivatives such a convenient disguise.
\end{keynote}

\section{Cut the domain and the difficulty disappears}

The next species does not collapse pointwise; it collapses after the interval is chopped. Absolute values, floors, fractional parts, signs, maxima and minima all define a partition of $[a,b]$ into pieces on which the integrand is a smooth, usually trivial, function. The work is entirely in finding the partition, and the failures are entirely off-by-one errors at its boundaries.

\tech[Absolute value, sign, and the modulus in the logarithm]{Absolute value, sign, and the log's modulus}{easy}
\cue $\abs{f(x)}$, $\sgn f(x)$, $\sqrt{x^{2}}$, $\sqrt{f(x)^{2}}$, or an odd-looking integrand on a symmetric interval where the bars have been inserted deliberately.
\method Bars \emph{change the parity}, and that is usually the whole point: $\abs{x}^{3}$ is even where $x^{3}$ is odd, so an integral you expected to vanish is instead $2\int_0^{a}$. Three rules cover the species. First, $\sqrt{f^{2}}=\abs{f}$, never $f$. Second, $\int\dfrac{\dd x}{x}=\ln\abs{x}$, and the modulus matters whenever the argument can be negative. Third, resolve the bars by splitting at every zero of the inner expression and choosing the sign on each piece.

\begin{worked}
\[\int_{-\infty}^{\infty}\abs{x}^{3}2^{-x^{2}}\dd x.\]
Without the bars this is odd and vanishes. With them the integrand is even, so
\[I=2\int_0^{\infty}x^{3}2^{-x^{2}}\dd x=2\cdot\frac12\int_0^{\infty}u\eu^{-u\ln2}\dd u=\frac{\Gfun(2)}{\ln^{2}2}=\frac{1}{\ln^{2}2}\approx2.081369,\]
substituting $u=x^{2}$ and then recognising a Gamma integral with rate $\ln2$. The bars rescued the integral from being $0$, and the entire content of the problem is noticing that.
\end{worked}

\begin{worked}[A radical that is secretly an absolute value, and a domain check]
\[\int_{-1}^{1}\sqrt{x^{13}+x^{8}}\dd x.\]
First the domain: $x^{13}+x^{8}=x^{8}(x^{5}+1)$, and on $[-1,1]$ both factors are nonnegative, so the radicand is nonnegative and the integrand is real throughout. Now $\sqrt{x^{8}}=\abs{x}^{4}=x^{4}$, because $8$ is even and the fourth power is nonnegative, so the integrand is $x^{4}\sqrt{1+x^{5}}$ with no bars needed. Substituting $u=1+x^{5}$, $\dd u=5x^{4}\dd x$:
\[I=\frac15\int_0^{2}\sqrt u\dd u=\frac{2}{15}\cdot2^{3/2}=\frac{4\sqrt2}{15}\approx0.377124.\]
Had the exponent been $x^{7}$ rather than $x^{8}$ the radicand would have been negative on part of $[-1,0)$ and the integral would not exist as a real object.
\end{worked}

\begin{trap}
Cancelling $\ln\abs{g}$ across a zero of $g$: the antiderivative $\ln\abs{g}$ is only valid on an interval where $g$ does not vanish, so an interior zero of $g$ makes the integral improper and the difference of logarithms meaningless. And $\sqrt{x^{2}}=x$ is the single most productive wrong step in this chapter: it converts $\int_0^{2\pi}\sqrt{1-\cos2x}\dd x=4\sqrt2$ into $0$, and the $0$ looks like a symmetry argument succeeding.
\end{trap}

\begin{seealsob}Splitting at a symmetry or sign-change point; even and odd parity on a symmetric interval; branch and domain traps; floor, fractional part, and step integrands.\end{seealsob}

Absolute values split the interval into two or three pieces. Floors and fractional parts split it into infinitely many, and the resulting infinite family is the point: the sum over pieces is a series, and the series is designed to telescope or to be a recognisable zeta value. The shells are usually easier to describe than the integrand.

\tech[Floor, fractional part, and step integrands]{Floor, fractional part, and step integrands}{medium}
\cue $\floorb{x}$, $\{x\}$, $\floorb{1/x}$, $\floorb{\sqrt{f(x)}}$, $\floorb{x^{1/3}}$, or any piecewise-constant weight multiplying something smooth.
\method Partition the domain into the level sets of the step function and integrate the now-smooth integrand on each. The three shell families worth memorising:
\[\floorb{x}=n\ \text{on}\ [n,n+1),\quad \floorb{\tfrac1x}=n\ \text{on}\ \Bigl(\tfrac{1}{n+1},\tfrac1n\Bigr],\ \text{length}\ \tfrac{1}{n(n+1)},\quad \floorb{x^{1/3}}=n\ \text{on}\ [n^{3},(n+1)^{3}),\ \text{length}\ 3n^{2}+3n+1.\]
Then sum the resulting series. The reciprocal family is the productive one, because $\tfrac{1}{n(n+1)}=\tfrac1n-\tfrac{1}{n+1}$ telescopes, and because it converts an integral on $(0,1]$ into a series indexed by shells that shrink quadratically.

\begin{worked}
\[\int_0^{1}\Bigl\{\frac1x\Bigr\}\dd x=1-\gamma\approx0.422784.\]
Substitute $u=\tfrac1x$: the integral becomes $\int_1^{\infty}\dfrac{\{u\}}{u^{2}}\dd u$. On $[k,k+1)$ we have $\{u\}=u-k$, so
\[\int_1^{N}\frac{\{u\}}{u^{2}}\dd u=\sum_{k=1}^{N-1}\int_k^{k+1}\Bigl(\frac1u-\frac{k}{u^{2}}\Bigr)\dd u=\ln N-\sum_{k=1}^{N-1}\frac{1}{k+1}\to1-\gamma,\]
using $\int_k^{k+1}\tfrac{k}{u^{2}}\dd u=\tfrac{1}{k+1}$ and the definition of $\gamma$ as $\lim\bigl(\sum_{1}^{N}\tfrac1k-\ln N\bigr)$. Direct shell-summing of the partial integrals confirms $0.4227818$.
\end{worked}

\begin{worked}[A cube-root shell, finite and exact]
\[\int_0^{64}\floorb{x^{1/3}}\dd x=\sum_{n=0}^{3}n\bigl((n+1)^{3}-n^{3}\bigr)=\sum_{n=0}^{3}n(3n^{2}+3n+1)=0+7+38+111=156.\]
The upper limit $64=4^{3}$ was chosen so that the last shell is complete; that is endpoint engineering in the service of a step function. Quadrature with break points at $1$, $8$, $27$ returns exactly $156$.
\end{worked}

\begin{trap}
Off-by-one at the shell boundaries. $\floorb{1/x}=n$ on the \emph{half-open} interval $\bigl(\tfrac{1}{n+1},\tfrac1n\bigr]$, and getting the direction wrong shifts the whole series by one index. The second error is an incomplete last shell: if the upper limit is not an integer (for $\floorb{x}$) or a perfect cube (for $\floorb{x^{1/3}}$), the final shell is truncated and contributes a fragment that must be computed separately. Third, check convergence before summing: $\int_0^{1}\floorb{1/x}\dd x=\sum_{n\geq1}\tfrac{n}{n(n+1)}=\sum_{n\geq1}\tfrac{1}{n+1}$ diverges, even though every shell is finite and the picture looks harmless.
\end{trap}

\begin{seealsob}Telescoping integrands and shell decomposition; absolute value, sign, and the log's modulus; periodicity and window shifting; splitting at a symmetry or sign-change point.\end{seealsob}

Shell decomposition produced a series that telescoped. That is not a coincidence of the floor function: it is the general design principle behind a whole family, in which either the domain decomposes and the resulting series telescopes, or the integrand itself is an infinite product or sum that collapses before any integration begins.

\tech{Telescoping integrands and shell decomposition}{hard}
\cue An infinite product such as $\prod_{n\geq0}(1+x^{2^{n}})$, an infinite sum of terms with $\tfrac{1}{n(n+1)}$-type weights, a difference $g(x+1)-g(x)$ inside the integrand, or a floor function partitioning the domain into shells on which the integrand repeats up to sign.
\method Two routes, and it is worth deciding which one you are on. If the \emph{integrand} telescopes, collapse it first: the binary identity
\[\prod_{n=0}^{\infty}\bigl(1+x^{2^{n}}\bigr)=\frac{1}{1-x}\quad(\abs{x}<1)\]
follows from $(1-x)\prod_{n=0}^{N}(1+x^{2^{n}})=1-x^{2^{N+1}}$. If the \emph{domain} decomposes, integrate on each shell and telescope the resulting series using $\tfrac{1}{n(n+1)}=\tfrac1n-\tfrac{1}{n+1}$ or a partial-fraction split of the shell weight.

\begin{worked}
\[\int_0^{1/2}\bigl(1+x\bigr)\bigl(1+x^{2}\bigr)\bigl(1+x^{4}\bigr)\bigl(1+x^{8}\bigr)\cdots\dd x=\int_0^{1/2}\frac{\dd x}{1-x}=\bigl[-\ln(1-x)\bigr]_0^{1/2}=\ln2\approx0.693147.\]
The product is a geometric series in disguise: multiplying by $(1-x)$ collapses it pairwise to $1-x^{2^{N+1}}$, which tends to $1$ on $[0,\tfrac12]$. A partial product with sixty factors integrates numerically to $0.6931471806$.
\end{worked}

\begin{worked}[A shell decomposition that telescopes]
\[\int_1^{\infty}\frac{\pi\sin(\pi x)}{\floorb{x}\bigl(\floorb{x}+1\bigr)}\dd x.\]
On $[n,n+1)$ the weight is the constant $\tfrac{1}{n(n+1)}$ and $\int_n^{n+1}\pi\sin(\pi x)\dd x=\bigl[-\cos(\pi x)\bigr]_n^{n+1}=2(-1)^{n}$. So
\[I=\sum_{n\geq1}\frac{2(-1)^{n}}{n(n+1)}=2\sum_{n\geq1}(-1)^{n}\Bigl(\frac1n-\frac{1}{n+1}\Bigr)=2\bigl[-\ln2-(\ln2-1)\bigr]=2-4\ln2\approx-0.772589.\]
The two alternating sums are $\sum(-1)^{n}/n=-\ln2$ and $\sum_{n\geq1}(-1)^{n}/(n+1)=\ln2-1$ after reindexing. Numerical shell-summing to $n=400$ gives $-0.7725950$, converging slowly as an alternating series must.
\end{worked}

\begin{trap}
Reindexing. Starting the binary product at $n=1$ instead of $n=0$ drops the factor $(1+x)$ and turns $\tfrac{1}{1-x}$ into $\tfrac{1}{1-x^{2}}$, a different answer entirely; and in the shell sum, starting at $n=0$ divides by zero. Write the first two shells out explicitly before summing. The second trap is convergence: a telescoping series whose terms do not tend to zero telescopes to nothing useful, and an infinite product converges only where it converges, so $\prod(1+x^{2^{n}})=\tfrac{1}{1-x}$ is false at $x=1$ and meaningless beyond.
\end{trap}

\begin{seealsob}Floor, fractional part, and step integrands; partial fractions with distinct linear factors; series truncation with a tail bound; recognising the emergent constant.\end{seealsob}

\section{Where the naive antiderivative is wrong}

Everything so far has been an opportunity. This section is a hazard. There is a species of problem in which an antiderivative exists, is correct as an antiderivative, and nevertheless gives the wrong number when evaluated at the two endpoints, because it is discontinuous somewhere inside the interval. The fundamental theorem requires $F$ to be continuous on $[a,b]$ and to satisfy $F'=f$ throughout; an $F$ that jumps has violated the first hypothesis, and no amount of correctness in the differentiation saves it.

Three constructions produce such an $F$ almost automatically, and between them they account for nearly every instance. The first is an $\arctan$ whose argument passes through $\pm\infty$: $\arctan(k\tan x)$ jumps by $\pi$ at every odd multiple of $\tfrac{\pi}{2}$, and integrands of the form $\dfrac{1}{a+b\sin^{2}x}$ produce exactly that antiderivative. The second is $\ln\abs{g}$ where $g$ has an interior zero: the antiderivative goes to $-\infty$ on one side and comes back from $-\infty$ on the other, and in fact the integral is improper there rather than merely discontinuous. The third is a substitution that is not injective, of which $t=\tan\tfrac x2$ across $x=\pi$ is the standard example. In all three cases the failure is silent, and in all three the same diagnostic catches it: evaluate the candidate $F$ at a dozen points across $[a,b]$ and look for a jump.

\tech{Branch and domain traps}{hard}
\cue A real cube root of a negative number; $\arcsec$; $\sqrt{x^{2}}$; an $\arctan$ whose argument runs to $\pm\infty$ inside the interval; $\operatorname{arcosh}$ evaluated below $1$; a Weierstrass substitution $t=\tan\tfrac x2$ on an interval containing $\pi$; or any exponent manipulation that changes the domain.
\method Establish the domain and the branch of every factor before manipulating, and check the candidate antiderivative for continuity on the closed interval. If $F$ jumps at an interior point $c$, integrate over $[a,c]$ and $[c,b]$ separately and add the one-sided limits:
\[\int_a^bf=\bigl[F\bigr]_a^{c^{-}}+\bigl[F\bigr]_{c^{+}}^{b},\]
which differs from $F(b)-F(a)$ by exactly the jump. For real odd roots, note that $x^{-1/3}$ preserves sign, so $\bigl(x^{-1/3}\bigr)^{-3}=x$ is genuinely odd on $[-a,a]$; the principal complex cube root does not preserve sign and would give a different answer.

\begin{worked}
\[\int_0^{\pi}\frac{\dd x}{1+3\sin^{2}x}.\]
Dividing through by $\cos^{2}x$ and substituting $t=\tan x$ gives the antiderivative $F(x)=\tfrac12\arctan(2\tan x)$, and $F'$ is the integrand wherever $\tan x$ is defined. But $F(\pi)-F(0)=0$, which is absurd for a strictly positive integrand. The cause is the jump at $x=\tfrac{\pi}{2}$, where $\tan x$ blows up and $F$ drops from $\tfrac{\pi}{4}$ to $-\tfrac{\pi}{4}$. Splitting there:
\[I=\Bigl[\tfrac12\arctan(2\tan x)\Bigr]_0^{\pi/2^{-}}+\Bigl[\tfrac12\arctan(2\tan x)\Bigr]_{\pi/2^{+}}^{\pi}=\frac{\pi}{4}+\frac{\pi}{4}=\frac{\pi}{2}\approx1.570796,\]
which quadrature confirms. The correct value exceeds the naive one by exactly the jump $\tfrac{\pi}{2}$.
\end{worked}

\begin{worked}[Two branch checks that go the other way]
\[\int_{-1}^{1}\bigl(x^{-1/3}\bigr)^{-3}\dd x.\]
Read as real functions, $x^{-1/3}$ is the real cube root of $\tfrac1x$ and preserves sign, so $\bigl(x^{-1/3}\bigr)^{-3}=x$ on $[-1,1]\setminus\{0\}$, the integrand is odd, and the integral is $0$. A computer algebra system using the principal complex branch will disagree, because for $x<0$ it returns a root with argument $\tfrac{\pi}{3}$ rather than $\pi$.

Second: $\dvop{x}\arcsec x=\dfrac{1}{\abs{x}\sqrt{x^{2}-1}}$, not $\dfrac{1}{x\sqrt{x^{2}-1}}$, because $\arcsec$ has range $[0,\pi]\setminus\{\tfrac{\pi}{2}\}$ and is increasing on both branches. Numerically at $x=-2$ the derivative is $+0.2886751$, which is $\tfrac{1}{2\sqrt3}$ and positive, confirming the modulus. Dropping it flips the sign of every integral over a negative interval.
\end{worked}

\begin{worked}[The substitution itself is discontinuous]
\[\int_0^{2\pi}\frac{\dd x}{2+\cos x}.\]
The Weierstrass substitution $t=\tan\tfrac x2$ turns this into $\int\dfrac{2\dd t}{3+t^{2}}$, and applying it with bounds $t=0$ to $t=\tan\pi=0$ returns $0$, again absurd for a positive integrand. The fault is not the antiderivative but the substitution: $t=\tan\tfrac x2$ has a pole at $x=\pi$, which is strictly inside $[0,2\pi]$, so it is not a bijection onto $\R$ there. Slide the window using periodicity to $[-\pi,\pi]$, where the pole sits at both endpoints and the map is a genuine bijection:
\[\int_{-\pi}^{\pi}\frac{\dd x}{2+\cos x}=\int_{-\infty}^{\infty}\frac{2\dd t}{3+t^{2}}=\frac{2\pi}{\sqrt3}\approx3.627599.\]
The lesson generalises: a substitution must be monotone and finite on the closed interval, and checking that costs one glance at where it blows up.
\end{worked}

\begin{trap}
Assuming the algebra system has thought about it. Symbolic systems default to principal complex branches for fractional powers and to formal antiderivatives that may jump, and they will disagree with the real-analysis answer in exactly the cases this technique is about. The practical protocol: plot or sample the candidate antiderivative at ten points across the interval, and if it jumps, split there. A cheaper version of the same check is the sign test used above, since a positive integrand cannot have a nonpositive integral.
\end{trap}

\begin{seealsob}Absolute value, sign, and the log's modulus; inverse and co-function collapse; Weierstrass substitution; endpoint limit discipline.\end{seealsob}

\section{A limit that only looks like an integral}

The last species is different in kind. Here the object is a limit of a family of integrals, the family concentrates on a set of measure zero, and the answer is read off the geometry of the concentration rather than obtained by integrating anything. Recognising the family is the whole technique, and the families are few.

It belongs in this chapter for the same reason as everything else here: the integral for any fixed $n$ is genuinely hard, and the limit is trivial, so a reader who evaluates first and takes the limit second will not finish. The structural cue is a parameter appearing in a specific pattern, namely a prefactor and a matching power inside a kernel that is a function of $n$ times something vanishing. Written abstractly, $K_{n}(u)=nK(nu)$ with $\int K=1$, which is the definition of an approximate identity. Once you see that shape the rest is arithmetic on the zeros of the inner function.

\tech{Nascent delta and limiting kernels}{hard}
\cue A parameter $n\to\infty$ multiplying a kernel that concentrates: $\dfrac{n}{1+n^{2}u^{2}}$, $\sqrt{\tfrac{n}{\pi}}\eu^{-nu^{2}}$, $\dfrac{\sin(nu)}{\pi u}$, $\dfrac{n}{2}\eu^{-n\abs{u}}$, composed with an inner function that has isolated zeros.
\method Identify the kernel and its total mass. As $n\to\infty$ the integral localises at the zeros of the inner function, and near a \emph{simple} zero $x_{0}$ of $g$ the substitution $u=g(x)$ is locally linear with $\dd u\approx\abs{g'(x_{0})}\dd x$, so the contribution is $\dfrac{\text{mass}}{\abs{g'(x_{0})}}$. Sum over the real zeros lying inside the interval. For the Lorentzian $\dfrac{n}{1+n^{2}u^{2}}$ the mass is $\pi$; for the Gaussian and the Fejer kernels it is $1$.

\begin{worked}
\[\lim_{n\to\infty}\int_{-\infty}^{\infty}\frac{n}{1+n^{2}(x^{4}-1)^{2}}\dd x.\]
The inner function $g(x)=x^{4}-1$ has real zeros $x=\pm1$, both simple, with $g'(x)=4x^{3}$ so $\abs{g'(\pm1)}=4$. The Lorentzian has mass $\pi$, so the limit is
\[\frac{\pi}{4}+\frac{\pi}{4}=\frac{\pi}{2}\approx1.570796.\]
At $n=10^{3}$ careful quadrature (with break points placed at the spikes) gives $1.571973$; at $n=10^{5}$, $1.570808$. The convergence is $O(n^{-1/2})$ in this instance, which is why a lazy quadrature at large $n$ misses the spikes entirely and returns nearly zero.
\end{worked}

\begin{worked}[A Gaussian kernel, and the same arithmetic]
\[\lim_{n\to\infty}\int_{-2}^{2}\sqrt{\frac{n}{\pi}}\,\eu^{-n(x^{2}-1)^{2}}\dd x.\]
The Gaussian family $\sqrt{n/\pi}\,\eu^{-nu^{2}}$ has total mass $1$, not $\pi$. The inner function again has simple zeros at $\pm1$ with $\abs{g'}=2$, so the limit is $\tfrac12+\tfrac12=1$. Quadrature with break points at the spikes gives $1.0018959$ at $n=100$ and $1.0000188$ at $n=10^{4}$, converging like $n^{-1}$. Notice that the only thing that changed from the Lorentzian case is the mass; the localisation arithmetic is identical, which is why it is worth memorising the masses ($\pi$ for $\tfrac{n}{1+n^{2}u^{2}}$, $1$ for the Gaussian, the Fejer and the Poisson kernels) rather than the formulas.
\end{worked}

\begin{trap}
Double roots. If $g'(x_{0})=0$ the local model is quadratic rather than linear, the scaling changes, and the contribution is not $\tfrac{\pi}{\abs{g'}}$; it can be infinite. Check that every zero in range is simple before applying the formula. Second, only zeros \emph{inside} the interval count, and a zero exactly at an endpoint contributes half. Third, complex zeros of $g$ contribute nothing: $x^{4}-1$ has four roots but only two of them are real, and using all four doubles the answer.
\end{trap}

\begin{seealsob}Laplace's method and Watson's lemma; bounding by envelopes; nascent delta as a distributional limit; named-function answers.\end{seealsob}

\section{Running the sweep}

The six tests of the sweep map onto the eight techniques of this chapter, and the mapping is worth having as a picture rather than a list, because in practice you run the tests in parallel: a single reading of the integrand should light up at most one of them.

\begin{center}
\begin{tikzpicture}[node distance=5mm and 6mm]
\node[dtest] (a) {Same value at two\\ sample points?};
\node[dleaf, left=10mm of a] (b) {Disguised constant:\\ prove the identity,\\ multiply by $b-a$};
\node[dtest, below=8mm of a] (c) {Steps, bars, floors,\\ $\max$ or $\min$ present?};
\node[dleaf, right=10mm of c] (d) {Split into level sets;\\ sum the shell series};
\node[dtest, below=8mm of c] (e) {A squared denominator,\\ or a power tower,\\ or $f^{2}{+}g^{2}$ downstairs?};
\node[dleaf, left=10mm of e] (f) {Perfect derivative:\\ guess $F$, differentiate,\\ compare};
\node[dleaf, below=8mm of e] (g) {Check branches and\\ interior jumps of $F$,\\ then integrate honestly};
\draw[dedge] (a) -- node[dlab,above]{yes} (b);
\draw[dedge] (a) -- node[dlab,right]{no} (c);
\draw[dedge] (c) -- node[dlab,above]{yes} (d);
\draw[dedge] (c) -- node[dlab,right]{no} (e);
\draw[dedge] (e) -- node[dlab,above]{yes} (f);
\draw[dedge] (e) -- node[dlab,right]{no} (g);
\end{tikzpicture}
\end{center}

The final leaf is not a technique but a warning, and it belongs at the bottom because it applies to every integral you ever compute, not only to the ones in this chapter. Before evaluating an antiderivative at two endpoints, confirm that it is continuous between them. That single check catches the $\tan\tfrac x2$ pole, the $\arctan(2\tan x)$ jump, the interior zero under a $\ln\abs{\cdot}$, and the branch cut of a fractional power, which between them account for most of the confidently wrong definite integrals in circulation.

\begin{keynote}[Why these are worth practising]
An integral that collapses is not a curiosity. It is a test of whether you read structure before reaching for method, and reading structure first is exactly the habit that pays on integrals that do not collapse. The reader who tests parity, samples the integrand, and checks the domain out of reflex will lose ten seconds on an honest problem and gain an hour on a rigged one. The reader who begins by choosing a technique will do the reverse. Every problem in this chapter was built by someone who knew which of those two readers they were writing for.
\end{keynote}

\section{Recognition matrix}
\begin{recog}{@{}p{0.30\textwidth}p{0.36\textwidth}p{0.28\textwidth}@{}}
\rowcolor{mist}\mh{If you see} & \mh{Reach for} & \mh{If that fails} \\
\midrule
\endhead
$x^{x}$, $x^{1/x}$, $u^{v}$ times a sum & Logarithmic differentiation: is it $F'$? & Check both endpoints for improperness \\
A stray $(1\pm\ln x)$ in a product & Guess $F=x^{x}$ or $x^{1/x}$ and differentiate & Parts, with $\ln x$ as the part \\
$\dfrac{\text{low degree}}{(\text{something})^{2}}$ & Reverse quotient rule: try $\dfrac{u}{v}$ & Partial fractions after all \\
$\dfrac{fg'-gf'}{f^{2}+g^{2}}$ & $\arctan\dfrac{g}{f}$; the Wronskian shape & Divide through by $f^{2}$, then $t=g/f$ \\
$f'g+fg'$ split across two terms & Reverse product rule: answer $\bigl[fg\bigr]_a^{b}$ & Parts once and compare \\
Mixed logarithm bases & Change of base; expect a constant & Sample the integrand twice \\
Nested or iterated radical & Fixed-point equation $y=\sqrt{x+y}$ or $y=\sqrt{xy}$ & Sample, then solve for $y$ \\
A determinant of antiderivatives & Evaluate the entries; expect $0$ & Row-reduce symbolically \\
$\arctan x+\arctan\tfrac1x$ & Constant $\tfrac{\pi}{2}$ on $(0,\infty)$ & Differentiate to confirm \\
Odd factor, even weight, $[-a,a]$ & Integral is $0$; check integrability & Recentre first if not symmetric \\
$\abs{x}^{\text{odd}}$ on $[-a,a]$ & Bars flip parity: $2\int_0^{a}$ & Split at $0$ and add \\
$\sqrt{f(x)^{2}}$ & It is $\abs{f}$; split at the zeros of $f$ & Never write $\sqrt{f^{2}}=f$ \\
$\sqrt{x^{13}+x^{8}}$-type radicand & Factor out the even power; check sign & Domain first, substitution second \\
$\floorb{1/x}$ on $(0,1]$ & Shells $\bigl(\tfrac{1}{n+1},\tfrac1n\bigr]$, length $\tfrac{1}{n(n+1)}$ & Telescope the series \\
$\{x\}$ or $\{1/x\}$ & $\{u\}=u-\floorb{u}$; expect $\gamma$ & Sum shells to $N$, then take $N\to\infty$ \\
$\floorb{x^{1/3}}$, $\floorb{\sqrt x}$ & Shells of length $3n^{2}{+}3n{+}1$, $2n{+}1$ & Watch for an incomplete last shell \\
$\prod_{n\geq0}(1+x^{2^{n}})$ & Telescopes to $\tfrac{1}{1-x}$ on $\abs{x}<1$ & Multiply by $(1-x)$ to see it \\
Shell weight $\tfrac{1}{n(n+1)}$ & Partial fractions, then telescope & Alternating: expect $\ln2$ \\
$\arctan(2\tan x)$ or $\tan\tfrac x2$ inside & Check for a jump at $\tfrac{\pi}{2}$ or $\pi$ & Split and add one-sided limits \\
$\arcsec$, $\operatorname{arcosh}$, odd roots & Range and sign conventions; keep $\abs{x}$ & Sample the derivative numerically \\
Positive integrand, nonpositive answer & An antiderivative jumped, or a pole is inside & Sign test before reporting \\
$\dfrac{n}{1+n^{2}g(x)^{2}}$ with $n\to\infty$ & Mass $\pi$ over $\abs{g'}$ at each real simple zero & Double root: rescale locally \\
\end{recog}
